\documentclass[11pt]{amsart}
\usepackage{amsmath,amssymb,amsthm,enumitem,hyperref}

\newtheorem{theorem}{Theorem}
\newtheorem{lemma}[theorem]{Lemma}
\newtheorem{corollary}[theorem]{Corollary}
\theoremstyle{definition}
\newtheorem{definition}[theorem]{Definition}
\newtheorem{remark}[theorem]{Remark}

\newcommand{\A}{\mathcal A}
\newcommand{\B}{\mathcal B}
\newcommand{\Fin}{\mathrm{Fin}}
\newcommand{\alg}{\mathrm{alg}}
\newcommand{\ult}{\mathrm{ult}}
\newcommand{\bc}{\mathfrak c}
\newcommand{\bb}{\mathfrak b}
\newcommand{\wh}{\widehat}

\title{Smooth Compactifications with Nonseparable Remainder under \(\mathfrak b=\mathfrak c\)}
\author{agent\_lane\_00}
\date{2026-07-03}

\begin{document}
\maketitle

\section*{Source and target}

Drygier and Plebanek call a compactification \(\gamma\omega\) \emph{smooth}
when the natural copy of \(c_0\) is complemented in \(C(\gamma\omega)\)
\cite{DP2016}.  For Boolean algebras
\(\Fin\subseteq\A\subseteq P(\omega)\), their Lemma 3.1 says that
\(K_\A\) is smooth if and only if there is a bounded sequence of finitely
additive measures \((\mu_n)\) on \(\A\) such that
\[
        \mu_n|\Fin=0,\qquad
        \mu_n(A)-\delta_n(A)\longrightarrow0
        \quad(A\in\A).
\]
The source proves under CH that a smooth compactification of \(\omega\) can
have nonseparable remainder.  It also states that the authors do not know the
ZFC answer and that their CH proof likely relaxes to \(\bb=\bc\).

This packet proves that advertised relaxation.

\section*{Theorem}

\begin{theorem}\label{thm:main}
Assume \(\bb=\bc\).  Then there is a smooth compactification
\(\gamma\omega\) of \(\omega\) such that
\(\gamma\omega\setminus\omega\) is nonseparable.
\end{theorem}

Equivalently, there is
\(\Fin\subseteq\A\subseteq P(\omega)\) such that \(\A/\Fin\) is not
\(\sigma\)-centered and \(K_\A\) is smooth.

\section*{Proof intuition}

The source CH proof runs an \(\omega_1\)-stage induction through countable
Boolean algebras.  Countability has two roles: the key density-killing step can
list old algebra elements, and CH lets the construction enumerate all countable
dense subsets of the initial Cantor-like remainder.

Under \(\bb=\bc\), \(\bc\) is regular because \(\bb\) is regular.  Thus at every
stage \(\xi<\bc\), the current algebra has size \(<\bb\).  The countable listing
in the source lemma can be replaced by the standard \(\bb\)-domination
finite-cover lemma used in \cite{DP2015}.

There remains one bookkeeping trap: a countable dense set in the final compactum
may have many possible extensions over intermediate algebras.  The successor
lemma below handles this by using two dense families.  Sets chosen from the
fixed initial dense family force every future extension of that initial family
to contain the new set \(X\).  Sets chosen from a dense family in the current
remainder give the outer-measure estimate needed to extend the smoothness
measures.

\section*{Auxiliary lemmas}

For \(\Fin\subseteq\B\subseteq P(\omega)\), write \(K_\B\) for the Stone
compactification and \(K_\B^*=K_\B\setminus\omega\).  If \(\nu\) is a
probability on \(\B\) and \(Z\subseteq\omega\), put
\[
        \nu_*(Z)=\sup\{\nu(B):B\in\B,\ B\subseteq Z\},\qquad
        \nu^*(Z)=\inf\{\nu(B):B\in\B,\ Z\subseteq B\}.
\]

\begin{lemma}[finite-cover lemma]\label{lem:finite-cover}
Let \(|\B|<\bb\).  For each \(i\in\omega\), let
\((S_i(k))_{k\in\omega}\) be a decreasing sequence in \(\B\) with empty
intersection.  Then there is \(g\in\omega^\omega\) such that, for
\[
        X=\bigcup_{i\in\omega}S_i(g(i)),
\]
every \(B\in\B\) with \(B\subseteq X\) is contained in a finite subunion
\(\bigcup_{i\le N}S_i(g(i))\).
\end{lemma}

\begin{proof}
This is the finite-cover lemma from \cite{DP2015}; we recall the short proof.
For finite \(m\) and \(\varphi\in\omega^m\), write
\[
        X_\varphi=\bigcup_{i<m}S_i(\varphi(i)).
\]
For fixed \(B\in\B\), the set of \(\varphi\in\omega^m\) such that
\(B\setminus X_\varphi\) is infinite has finitely many minimal elements under
coordinatewise order, by induction on \(m\).  Hence there is a finite
\(I_m^B\subseteq\omega\) meeting \(B\setminus X_\varphi\) whenever that set is
infinite.

Choose \(h_B\in\omega^\omega\) so that
\[
        S_i(h_B(i))\cap\bigcup_{m\le i}I_m^B=\emptyset
        \qquad(i\in\omega).
\]
Since \(|\B|<\bb\), there is \(g\) eventually dominating all \(h_B\).  If
\(B\subseteq X\) and \(N\) is such that \(h_B(i)\le g(i)\) for \(i\ge N\), then
\(B\setminus X_{g|N}\) cannot be infinite, because any point in
\((B\setminus X_{g|N})\cap I_N^B\) would have to lie in a tail
\(S_i(g(i))\), \(i\ge N\), disjoint from \(I_N^B\).  Thus
\(B\setminus X_{g|N}\) is finite, and increasing \(N\) captures those finitely
many points.
\end{proof}

\begin{lemma}[successor step]\label{lem:successor}
Let \(\Fin\subseteq\A_0\subseteq\B\subseteq P(\omega)\), with \(\A_0\)
countable and \(|\B|<\bb\).  Let \(D=\{d_j:j\in\omega\}\) be a countable dense
subset of \(K_{\A_0}^*\), and let \(Q=\{q_j:j\in\omega\}\) be a countable dense
subset of \(K_\B^*\).  Suppose that \((\nu_n)\) is a sequence of nonatomic
probabilities on \(\B\) such that
\[
        \nu_n|\Fin=0,\qquad
        \nu_n-\delta_n\to0
        \quad\hbox{pointwise on }\B .
\]
Then there is an infinite \(X\subseteq\omega\) such that:
\begin{enumerate}[label=(\roman*)]
\item \(\omega\setminus X\) is infinite;
\item every ultrafilter on \(\B[X]\) whose restriction to \(\A_0\) is one of
the \(d_j\)'s contains \(X\);
\item the \(\nu_n\)'s extend to nonatomic probabilities
\(\widetilde\nu_n\) on \(\B[X]\) with
\(\widetilde\nu_n-\delta_n\to0\) pointwise on \(\B[X]\).
\end{enumerate}
\end{lemma}

\begin{proof}
Enumerate the disjoint union \(D\sqcup Q\) as \((r_i)_{i\in\omega}\), retaining
the information whether \(r_i\) is an ultrafilter on \(\A_0\) or on \(\B\).
Choose positive numbers \(\eta_i\) with \(\sum_i\eta_i<1/2\).

For every \(i\) we build a decreasing sequence \(S_i(k)\), \(k\in\omega\), with
empty intersection.  If \(r_i=d_j\in K_{\A_0}^*\), then
\[
        S_i(k)\in\A_0\cap d_j .
\]
If \(r_i=q_j\in K_\B^*\), then
\[
        S_i(k)\in\B\cap q_j .
\]
In both cases we also require
\[
        \nu_m(S_i(k))<\eta_i\quad(m\le k),                 \tag{1}
\]
and
\[
        n\notin S_i(k)\quad\Longrightarrow\quad
        \nu_n(S_i(k))<\eta_i .                              \tag{2}
\]

Assume \(S_i(k-1)\) has been chosen, with \(S_i(-1)=\omega\).  Work in
\(\A_0\) if \(r_i\in D\), and in \(\B\) if \(r_i\in Q\).  Since the finite sum
\(\nu_0+\cdots+\nu_k\), restricted to that algebra, is nonatomic and \(r_i\) is
nonprincipal, choose
\[
        C\in r_i,\qquad
        C\subseteq S_i(k-1)\setminus\{0,\ldots,k-1\},
\]
with \(\nu_m(C)<\eta_i\) for \(m\le k\).  By pointwise convergence on \(C\),
the set
\[
        F=\{n\notin C:\nu_n(C)\ge\eta_i\}
\]
is finite.  The previous step implies \(F\subseteq S_i(k-1)\), and the choice
of \(C\) gives \(F\cap\{0,\ldots,k-1\}=\emptyset\).  Put
\[
        S_i(k)=C\cup F .
\]
Finite sets have \(\nu_n\)-measure zero, so (1) and (2) hold.  The sequence is
decreasing and eventually avoids every integer.

Apply Lemma \ref{lem:finite-cover} to \((S_i(k))\), and set
\[
        X=\bigcup_i S_i(g(i)).
\]
If \(\omega\subseteq X\), the finite-cover property gives
\(\omega\subseteq\bigcup_{i\le N}S_i(g(i))\) for some \(N\), and hence
\[
        1=\nu_0(\omega)\le \sum_{i\le N}\nu_0(S_i(g(i)))
        \le \sum_{i\le N}\eta_i<1/2,
\]
impossible.  Thus \(\omega\setminus X\) is infinite.

For every \(d_j\in D\), one of the chosen sets \(S_i(g(i))\) lies in \(d_j\)
and is contained in \(X\).  Therefore every ultrafilter on \(\B[X]\) extending
\(d_j\) contains \(X\).  This proves (ii).

It remains to extend the measures.  Fix \(n\).  If \(B\in\B\) and
\(\nu_n(B)>0\), then \(B\) is infinite.  The clopen \(\wh B\cap K_\B^*\) is
nonempty, so the dense set \(Q\) contains some \(q_j\) with \(B\in q_j\).
For the corresponding index \(i\), both \(B\) and \(S_i(g(i))\) lie in
\(q_j\), so \(B\cap X\ne\emptyset\).  Hence no positive-\(\nu_n\) old algebra
element is disjoint from \(X\), and
\[
        \nu_n^*(X)=1 .                                      \tag{3}
\]

Now let \(n\notin X\).  If \(B\in\B\) and \(B\subseteq X\), then
\[
        B\subseteq\bigcup_{i\le N}S_i(g(i))
\]
for some \(N\).  For fixed \(K\le N\), split the right side into the first
\(K\) pieces and the tail.  By (2),
\[
        \nu_n\!\left(\bigcup_{K\le i\le N}S_i(g(i))\right)
        \le \sum_{i\ge K}\eta_i .
\]
For each \(i<K\), pointwise convergence on \(S_i(g(i))\) gives
\(\nu_n(S_i(g(i)))\to0\) along \(n\notin X\).  Thus
\[
        \limsup_{n\notin X}(\nu_n)_*(X)\le \sum_{i\ge K}\eta_i,
\]
and letting \(K\to\infty\) gives
\[
        \lim_{n\notin X}(\nu_n)_*(X)=0.                     \tag{4}
\]

Define \(t_n=1\) for \(n\in X\) and \(t_n=(\nu_n)_*(X)\) for \(n\notin X\).
By (3), \(t_n\in[(\nu_n)_*(X),\nu_n^*(X)]\).  The Los--Marczewski extension
theorem gives a probability extension \(\widetilde\nu_n\) to \(\B[X]\) with
\(\widetilde\nu_n(X)=t_n\).  By (4),
\(\widetilde\nu_n(X)-\delta_n(X)\to0\), so the source convergence lemma
\cite[Lemma 3.3]{DP2016} yields
\[
        \widetilde\nu_n-\delta_n\to0
        \quad\hbox{on }\B[X].
\]

The extensions are nonatomic.  If \(u\in\ult(\B[X])\) is principal, then a
singleton in \(u\) has measure zero.  Otherwise \(u|\B\) is nonprincipal, and
nonatomicity of \(\nu_n\) on \(\B\) gives \(B\in u|\B\) with
\(\nu_n(B)<\varepsilon\).  One of \(B\cap X\) and \(B\setminus X\) lies in
\(u\), and its \(\widetilde\nu_n\)-measure is at most \(\nu_n(B)\).
\end{proof}

\section*{Proof of Theorem \ref{thm:main}}

Let \(\A_0\) and \((\nu_n^0)\) be as in Lemma 6.2 of \cite{DP2016}: \(\A_0\)
is countable, \(\Fin\subseteq\A_0\), each \(\nu_n^0\) is a nonatomic
probability vanishing on finite sets, and
\[
        \nu_n^0-\delta_n\to0
        \quad\hbox{on }\A_0 .
\]
The space \(K_{\A_0}^*\) is compact metrizable, so it has exactly \(\bc\) many
countable dense subsets.  Enumerate them as
\[
        (D_\xi)_{\xi<\bc}.
\]

We construct an increasing continuous chain
\((\A_\xi,\nu_n^\xi)_{\xi<\bc}\).  At limit stages take unions.  Since
\(\bb=\bc\) and \(\bb\) is regular, every \(\A_\xi\) has cardinality
\(<\bb\).

At stage \(\xi\), suppose \(\A_\xi\) and \((\nu_n^\xi)\) have been constructed.
If \(K_{\A_\xi}^*\) is nonseparable, stop.  Then \(K_{\A_\xi}\) is already a
smooth compactification by the measure criterion, and its remainder is
nonseparable.

Otherwise choose a countable dense set \(Q_\xi\subseteq K_{\A_\xi}^*\).  Apply
Lemma \ref{lem:successor} with
\[
        \B=\A_\xi,\qquad D=D_\xi,\qquad Q=Q_\xi,
        \qquad \nu_n=\nu_n^\xi .
\]
Let \(X_\xi\) be the set produced by the lemma, put
\[
        \A_{\xi+1}=\A_\xi[X_\xi],
\]
and let \(\nu_n^{\xi+1}\) be the corresponding extensions.  At a limit
\(\lambda<\bc\), set
\[
        \A_\lambda=\bigcup_{\xi<\lambda}\A_\xi,
\]
and let \(\nu_n^\lambda\) be the unique union of the coherent earlier measures.
These limit measures are nonatomic because their restrictions to the fixed
initial algebra \(\A_0\) are nonatomic.

If the construction stops, we are done.  Suppose it never stops, and define
\[
        \A=\bigcup_{\xi<\bc}\A_\xi,\qquad
        \mu_n=\bigcup_{\xi<\bc}\nu_n^\xi .
\]
For every \(A\in\A\), the set \(A\) belongs to some stage, so
\[
        \mu_n(A)-\delta_n(A)\to0.
\]
Thus \(K_\A\) is smooth.

We prove \(K_\A^*\) is nonseparable.  Suppose \(E=\{e_j:j\in\omega\}\) were
dense in \(K_\A^*\).  The restriction map \(K_\A^*\to K_{\A_0}^*\) is onto, so
\[
        D=\{e_j|\A_0:j\in\omega\}
\]
is dense in \(K_{\A_0}^*\).  Hence \(D=D_\xi\) for some \(\xi<\bc\), after
renumbering.  At stage \(\xi\), Lemma \ref{lem:successor} made every future
extension of a member of \(D_\xi\) contain \(X_\xi\).  Therefore every
\(e_j\) contains \(X_\xi\).

But \(\omega\setminus X_\xi\) is infinite and belongs to \(\A_{\xi+1}\), hence
the clopen set
\[
        \wh{\omega\setminus X_\xi}\cap K_\A^*
\]
is nonempty.  It is disjoint from \(E\), contradicting the density of \(E\).
So \(K_\A^*\) is nonseparable, and Theorem \ref{thm:main} follows.

\section*{Novelty and limitations}

The run indexes had no previous full packet for arXiv:1601.03770.  A bounded
web search for the exact source title, ``smooth compactification'',
``nonseparable remainder'', and ``\(\bb=\bc\)'' found the source paper
\cite{DP2016} and the companion nonseparable-growth paper \cite{DP2015}, but no
later theorem proving the smooth \(\bb=\bc\) strengthening.

This packet proves the \(\bb=\bc\) version suggested in the source.  It does
not prove the stronger ZFC assertion.

\section*{Human review recommendation}

First review Lemma \ref{lem:successor}.  In particular, check that the
\(\A_0\)-pieces force all future extensions of \(D_\xi\) to contain \(X_\xi\),
while the current dense set \(Q_\xi\) is sufficient for \(\nu_n^*(X_\xi)=1\).
Then review the final nonseparability contradiction.

\begin{thebibliography}{9}

\bibitem{DP2016}
P. Drygier and G. Plebanek,
\emph{Compactifications of \(\omega\) and the Banach space \(c_0\)},
arXiv:1601.03770, 2016.

\bibitem{DP2015}
P. Drygier and G. Plebanek,
\emph{Nonseparable growth of the integers supporting a measure},
arXiv:1501.06972, 2015.

\end{thebibliography}

\end{document}
